\section{Disk-certified transfer after an adaptive pause}
\label{sec:disk-transfer-integration}

Report 21 supplies the geometric certificate needed to apply a finite radical
transfer to actual scan prefixes. The integrated option
\code{--reduction disk-adaptive} extends the existing adaptive policy:
ordinary sparse cancellation first, the cheap binary survivor test second,
and a bounded full-transfer attempt only if adjacent surviving degrees prevent
that shortcut. Successful transfer retains the radical differential between
those degrees. Failure of the geometric or local-work gate resumes the same
partially reduced complex. The default remains \code{standard}.
This is an implementation of another exact cancellation strategy within the
local-complex framework of Bar-Natan \cite{barnatan}, not a new recognition
criterion or a general quasi-polynomial algorithm.

\subsection{Certifying the boundary from the processed rotation system}
For a processed prefix $P$, pair the darts of every internal edge and fix each
frontier dart. Call this permutation $\alpha$; let $\sigma$ rotate the four
darts at each crossing. The cycles of $\sigma\alpha$ are boundary walks of
the ribbon neighborhood of $P$. The implementation checks that this prefix
is connected, that
\[
 |V(P)|-|E_{\rm internal}(P)|+\#\text{boundary walks}=2,
\]
and that all frontier darts occur on one boundary walk. These conditions
certify a genus-zero neighborhood with one marked boundary component.
Filling its unmarked holes gives a disk. In a validated connected full
diagram, no unprocessed crossing component can hide in such a hole: a
connection to the rest would produce a frontier dart on that boundary.

Reading the marked boundary gives the common cyclic endpoint order. Every
smoothing is embedded in this disk, so its matching is noncrossing in that
order. Delooping closed circles leaves endpoint pairings unchanged. The code
also checks all distinct live matching types against the order, using the
usual stack test for nested arcs. Numeric edge labels do not themselves
specify an order. In particular the three pairings of four points cannot all
be noncrossing in one cyclic order; applying a disk-algebra theorem to them
without this check would be invalid.

The certificate records the boundary walks, vertex and internal-edge counts,
cyclic order, and a digest of the actual prefix. The digest binds data but
does not authenticate a source. Replay recomputes all the fields. Certificate
construction is linear in indexed prefix darts; checking every prefix
independently costs $O(n^2)$ such graph operations. The adapter records the
actual supplied scan order and verifies its current frontier against the
certificate. It does not replace a good greedy order automatically.

The report also supplies a polynomial ear-insertion construction of connected
prefix-and-suffix orders for biconnected crossing graphs. Such a cut admits
the disk certificate, but the construction need not preserve a small width.
The helper is retained for experiments and independently checked on smoothing
states; no claim of a width approximation or a universally good order follows.

\subsection{Retaining the maps that the binary profile cannot determine}
For a frontier of $2b$ points, let $J$ be the positive-weight radical in its
disk matching category. A dotted basis map from matching $a$ to matching $c$
has weight $b-k(a,c)+2|S|$, where $k(a,c)$ is the number of overlay circles.
Composition adds weights; the only weight-zero basis maps are identity maps
of equal matchings. The earlier radical analysis gives $J^{2b}=0$ for
$b\ge1$. Briefly, the weight bound first gives $J^{2b+1}=0$; a potentially
nonzero product of $2b$ positive factors would consist entirely of undotted
weight-one surfaces and produce a fully dotted endomorphism. Over $\F$, an
undotted positive-genus component vanishes and every surviving genus-zero
component leaves an output circle undotted. The fully dotted term is
therefore impossible. If all positive maps preserve their matching, the
square-free endomorphism ring gives the stronger exponent $b+1$.

Split the actual differential as $d=d_0+\delta$, retaining in $d_0$ only
the constant coefficient of maps between equal matchings. Undotted maps
between distinct matchings belong to $\delta$; bit zero alone is not an
identity test. Binary elimination constructs a contraction onto the residue
homology with inclusion $i$, projection $p$, and degree-minus-one homotopy $h$:
\[
 pi=1,\quad d_0h+hd_0=1+ip,\quad hi=ph=h^2=0,\quad d_0i=pd_0=0.
\]
All three maps have scalar identity coefficients. Hence nilpotence makes
the following inverse series finite:
\begin{align*}
 I&=(1+h\delta)^{-1}i,&
 D&=p\delta(1+h\delta)^{-1}i,\\
 H'&=(1+h\delta)^{-1}h,&
 P&=p(1+\delta h)^{-1}.
\end{align*}
These are the characteristic-two perturbation formulas; the standard
homological perturbation lemma supplies the deformation retract
\cite{crainic-perturbation}. In particular
\[
 D=\sum_{j=0}^{2b-2}p\delta(h\delta)^ji
\]
for $b\ge1$, with at most $b$ positive factors needed in the
matching-preserving case. At a closed frontier $D=0$ and binary homology
already suffices. The output has exactly the previously predicted minimal
matching multiplicities, with its nonzero adjacent-degree maps intact.

The kernel streams one residual column at a time: apply $\delta$, accumulate
its projection by $p$, and replace the current vector by its image under
$h\delta$. It computes an additional final zero check rather than silently
cutting off a nonzero term. Unexpected failure of nilpotence raises an
algebraic error; it does not become a local resource decline or a knot verdict.
Optional full certificates include $I,P,H'$ and are verified by separate
typed matrix multiplication. The checks include both square-zero equations,
both chain-map equations, $PI=1$, $dH'+H'd=1+IP$, and the three side
conditions $H'I=PH'=(H')^2=0$.

\subsection{The adaptive implementation and its transactional boundary}
As before, a stage permits $\max(256,4N)$ ordinary Schur-update pairs before
pausing between complete pivots. A successful cheap residue/degree-gap test
finishes the stage immediately. Otherwise the disk option takes a compact
snapshot and attempts the geometric certificate and finite transfer.
The new differential and reverse incidences are built separately, and all
five object arrays are installed together only after the last budget check.
The old coefficient-result cache is cleared at installation; typed algebra
plans remain valid. A declined attempt leaves the old differential and its
completed sparse cancellations intact.

The default experimental allowance is $\max(4096,64N)$ cooperative polls for
one attempt. Geometry traversal, contraction columns, series steps, coefficient
products, and scalar incidences poll this guard. The global deadline is
checked first at each poll; global exhaustion still raises \code{ScanLimit}
and becomes \code{UNKNOWN} in recognition. Local exhaustion only declines
the attempt. This poll allowance is neither a bit-operation budget nor a
memory ceiling: binary vector work and a coefficient composition can cost
more than one elementary operation. There is no competitive-runtime theorem
for this scheduler. Direct kernel calls can perform the complete finite
transfer without that experimental allowance.

The new policy is available to raw homology, recognition, allocation-pruned
windows, and minimal windows. The window adapter records each processed
crossing before either its ordinary or pruned extension calls cancellation,
so the geometric certificate remains bound to the correct prefix. Existing
component coefficient engines are compatible; component-owner arrays,
saturated weights, races, and alternative scanner representations are not
silently reinterpreted. The same standard-scanner restrictions as the earlier
adaptive policy apply. Evidence includes attempts, completed transfers,
geometry declines, local-budget fallbacks, transfer compositions, and maximum
radical depth. Full replay certificates are an explicit kernel diagnostic,
not an extra default allocation during recognition.

\subsection{Cost separation and the missing global parameter bound}
Let $N$ be the pre-reduction object count, $r$ the residue survivor count,
$E_\delta$ the number of nonzero radical entries, and $S_h,S_p$ the numbers
of scalar incidences in the homotopy and projection. Write $A_b,M_b$ for
coefficient addition and correctly typed composition costs. The column
algorithm has the conservative mixed-cost bound
\[
 O(N^3)+O\!\left(N+E_\delta+
 \ell r\bigl(E_\delta M_b+(S_h+S_p)A_b\bigr)\right),\quad \ell\le2b.
\]
The first term is binary contraction. Dense coefficient work is thus
$O(brN^2M_b)$, instead of a generic $O(N^3M_b)$ scalar-pivot bound. It can
help when $r\ll N$ and sparse elimination would create expensive fill.
It need not help sparse matrices. The product cache is bounded at 4096
entries; optional full certificates require additional storage and work.

For a fixed disk-certified full scan, let $B$ be its maximum boundary
half-width and $R$ its largest minimal explicit object count. Attaching one
crossing creates at most eight objects per old object, so stage input size
is at most $8R$. Coefficients occupy at most $2^B$ bits and even elementary
composition has cost $\operatorname{poly}(B)2^{O(B)}$. Noncrossing matching
types number at most $4^B$. Charging algebra caches and geometry as well as
matrices yields
\[
 \operatorname{poly}(n,1+R)2^{O(B)}.
\]
This conditional explicit-scan bound is compatible with ordinary exhaustive
pivots too; the new transfer improves how expensive coefficient work is
organized. It does not prove smaller intrinsic minimal multiplicities.
The condition $B+\log(1+R)=O(\log^2 n)$ is sufficient for quasi-polynomial
time, but no universal extraction theorem for those parameters is supplied.

\subsection{Why changing crossing order alone cannot suffice}
Report 21 constructs an unknot diagram on an $m\times m$ crossing grid,
with $n=m^2$. Pair exterior endpoints adjacent in their remaining cyclic
order, always joining distinct open paths until only one closed strand
remains. This adds no crossings and preserves the internal grid as a
spanning subgraph. Make the first visit to each crossing an overpass; the
result is descending and hence an unknot.

Every ordering has a prefix containing between $m^2/3$ and $2m^2/3$
vertices. If at least $m/3$ rows are mixed across the prefix cut, each
contributes a horizontal cut edge. Otherwise there is both a full and an
empty row: mixed rows alone cannot contain either side's required size.
Each column then contributes a vertical cut edge between those rows.
Thus every ordering has at least $m/3$ raw frontier edges, or
$B\ge m/6=\Omega(\sqrt n)$.
This refutes a universal polylogarithmic \emph{raw-order-width} strategy.
It is not a recognition lower bound: the descending certificate and ordinary
Reidemeister simplification bypass these diagrams. In the production audit,
all generated inputs through 64 crossings were recognized before homology;
independent cube ranks were checked only through nine crossings.

\subsection{Validation and measured scope}
The report's 52 tests pass unchanged. Its source-derived fixture was audited
against the pinned Git blobs for \code{scan\_fast.py}, \code{planar.py}, and
the scoped geometry definitions; executable ASTs agree after stripping
documentation. The production suite now passes 344 tests, including 100
seeded small diagram comparisons with an independent cube, intermediate
minimal-profile comparisons, adjacent-map synthetic checks, failure atomicity,
both window implementations, coefficient composition, and global limits.
An eager test control deliberately exercises transfer on actual prefixes;
this is distinct from observing a production adaptive switch.

The separately rerun geometric audit accepted 71 of 80 order candidates,
checked 14,658 complete smoothing states and 71 independent cube ranks,
and generated and verified seven typed full chain-homotopy certificates.
Data and the source audit are retained under \path{synthesis/data/}; the
reproduction driver is \path{fast/check_disk_geometry.py}.

Seven shuffled paired rounds in \path{fast/benchmark_disk_transfer.py}
compare fresh states under ordinary, A/A, residue-adaptive, and disk-adaptive
policies. Dense synthetic two-term blocks with nonzero survivor maps gained
1.126, 2.948, 4.078, and 3.307 times for the four retained cases; every disk
run made one successful full-transfer call. Independent binary regular-
representation homology agrees before and after reduction. These synthetic
blocks, including their coefficient-degree mixtures, are not asserted to be
actual knot-prefix complexes.

The sparse diagonal control stayed on ordinary pivots: no transfer attempt
was made, and its ratio was 0.959 rather than the eager prototype's large
regression. The four full diagram scans (Conway, Kinoshita--Terasaka, the
36-crossing stress braid, and the eight-crossing hard unknot) made zero full
transfer calls. Their ratios ranged from 0.960 to 1.010, while A/A controls
ranged from 0.992 to 1.012. Normal recognition decided all four in earlier
certificate or invariant stages. These measurements support the overhead
guard on these controls; they establish no recognition speedup from full
transfer. Defaults remain unchanged and the goal of general quasi-polynomial
recognition remains open.
